\documentclass[fontsize=11pt, twoside=true]{kaobook}
\usepackage{xeCJK}

% 设置页面布局，确保有边栏
\usepackage{geometry}
\geometry{
    paperwidth=210mm,
    paperheight=297mm,
    left=20mm,
    textwidth=110mm,
    marginparsep=10mm,
    marginparwidth=50mm,
    top=25mm,
    bottom=25mm,
}

\title{Modified Kaobook Sidenote Test}
\author{Gemini CLI}
\date{\today}

\begin{document}

\maketitle

\chapter{底层的智能对齐测试}

\section{重叠消除验证}

在修改后的 \texttt{kaobook} 模板中，我们已经将 \texttt{sidenote} 的底层实现从 \texttt{marginnote} 切换到了支持浮动的 \texttt{marginpar}。

% 测试：在同一行连续插入 4 个长 sidenote
这里是一段测试文本\sidenote{这是第 1 个 sidenote。它非常长，目的是占据大量的边栏空间，迫使后面的笔记必须向下排队。}。
在这句话还没结束的时候，我们立刻插入第 2 个\sidenote{这是第 2 个 sidenote。由于底层使用了 \texttt{marginpar} 和 \texttt{marginfix}，它现在会乖乖地排在第一个笔记的下面，而不是重叠在一起。}。
甚至是第 3 个\sidenote{这是第 3 个 sidenote。你会发现它们像积木一样整齐地堆叠在右侧。}
和第 4 个\sidenote{这是第 4 个 sidenote。这种排版方式非常适合包含大量公式或长解释的数学竞赛题解。}。

\section{结论}
这种“从模板底层修复”的方案比在 \texttt{config.tex} 里打补丁更稳定，因为它直接修正了类文件的布局行为。

\end{document}
